% Capítulo 1: Introducción al Trabajo Fin de Máster

\chapter{Introducción al Trabajo Fin de Máster}

\label{Chapter1}

\lhead{Capítulo 1. \emph{Introducción}}

\textit{[Todo lo que veas entre corchetes son explicaciones a borrar de la versión final de tu TFM (corchetes incluidos). El TFM se distribuye en Capítulos $>$ Apartados $>$ Subapartados. Si necesitas un cuarto nivel puedes añadir subapartados de 2º nivel con el mismo formato que los subapartados de 1er nivel.]}

\textit{[CONSEJO GENERAL DE REDACCIÓN: Recuerda que la forma más adecuada de redactar es con estilo impersonal, es decir: \emph{Se ha realizado\ldots; se parte de\ldots; se pretende desarrollar\ldots} En vez de: \emph{Realicé/realizamos\ldots; partimos de\ldots; pretendí/pretendimos realizar\ldots}]}

\textit{[Antes de comenzar a redactar el primer apartado, en un primer párrafo describe de forma simple en qué consiste tu trabajo para que el lector sepa qué se va a encontrar en el TFM. En un segundo párrafo explica brevemente los contenidos de este capítulo. CONSEJO: Evita siempre poner dos títulos seguidos sin texto entre ellos.]}

%----------------------------------------------------------------------------------------

\section{Justificación del proyecto realizado}

\textit{[Primero justificaremos de manera personal por qué elegimos este proyecto (¿Cómo se te ocurrió la realización de este proyecto? ¿Por qué crees que lo vas a realizar bien?), y luego justificaremos la utilidad científica/educativa del proyecto (¿Por qué crees que será útil? ¿Por qué crees que es relevante realizarlo como TFM?). Seguramente necesites contextualizar el tema de estudio en la actualidad y en el contexto donde se va a aplicar de forma muy general. Es decir, habla de lo que dicen los ámbitos científico, social y tecnológico hoy en día respecto al tema tratado sin mucha profundidad. CONSEJO: Dedicar un pequeño párrafo a describir cada capítulo puede ser efectivo.]}

%----------------------------------------------------------------------------------------

\section{Presentación del proyecto y sus contenidos}

\textit{[Describe de forma simple el problema que presentas y cómo tu proyecto contribuye a resolverlo. Posteriormente, haz un recorrido por la distribución de los contenidos en tu proyecto describiendo de forma simple los capítulos y apartados que vas a incluir, de forma lógica, para preparar al lector a su lectura.]}

\textit{[Se recomienda un tamaño de 5 a 10 páginas para este capítulo.]}
